\documentclass[12pt]{article}
\usepackage[utf8]{inputenc}
\usepackage[T1]{fontenc}
\usepackage{hyperref}
\title{Yuwei Language Resources for K-12 Curriculum Development}
\author{Compiled by Claude Sonnet 5.5 (claude-sonnet-5-5) for Andr\'as Kornai}
\date{}
\begin{document}
\maketitle

\section{Introduction}
The code whg is North Wahgi, also called Yu We or Yuwei, of Western Highlands Province, Papua New Guinea. It is a separate code from Mid-Wahgi (wgi). It should not be confused with Western Province languages such as Were or Yonggom, or with Muyuw and Yaweyuha, which also came up in searches. Joshua Project says mother-tongue literacy is very low and that no active literacy programmes exist.

\section{Primary Educational Resources}
No primer, textbook or school programme was found.

\section{Linguistic and Reference Materials}
\begin{itemize}
\item Heather and Don McLean (2005), \emph{North Wahgi (Yu We) language [Western Highlands Province]}, a phonological description, listed in Glottolog: \url{https://glottolog.org/resource/reference/id/109907}
\item Wikipedia, Wahgi language (the related Mid-Wahgi, wgi): \url{https://en.wikipedia.org/wiki/Wahgi_language}
\end{itemize}
No dictionary was found.

\section{Religious and Community Resources}
\begin{itemize}
\item Joshua Project, North Wahgi: \url{https://joshuaproject.net/people_groups/15819}
\item Scripture Earth listing for whg: \url{https://www.scriptureearth.org/00eng.php?iso=whg}. I could not confirm how much of the Bible is published.
\end{itemize}

\section{Digital and Audio Resources}
\begin{itemize}
\item Global Recordings Network, Yuwei (whg): \url{https://globalrecordings.net/en/language/whg}
\item Global Recordings Network, North Wahgi (Banz-Nondugl): \url{https://globalrecordings.net/en/language/5166}
\end{itemize}
The audio is evangelistic and Bible-teaching material, not language-learning content.

\section{Curriculum Development Context}
The language is not known to be taught in schools. The available material is a phonological description, scripture-related items and audio recordings. These could support an orthography and a first reader, but a curriculum would need new community-led material.
\end{document}